% Generated draft: requires mathematical and source review.
\documentclass[11pt]{article}
\usepackage[margin=25mm]{geometry}
\usepackage{fontspec}
\setmainfont{Latin Modern Roman}
\usepackage{amsmath,amssymb,mathtools,unicode-math}
\setmathfont{Latin Modern Math}
\usepackage{enumitem,needspace,xurl,hyperref}
\setlength{\parindent}{0pt}
\setlength{\parskip}{5pt}
\setlength{\emergencystretch}{3em}
\newcommand{\N}{\mathbb{N}}
\newcommand{\Z}{\mathbb{Z}}
\newcommand{\Q}{\mathbb{Q}}
\newcommand{\R}{\mathbb{R}}
\newcommand{\C}{\mathbb{C}}
\newcommand{\F}{\mathbb{F}}
\newcommand{\A}{\mathbb{A}}
\newcommand{\PP}{\mathbb{P}}
\newcommand{\E}{\mathbb{E}}
\newcommand{\eps}{\varepsilon}
\newcommand{\dd}{\,\mathrm{d}}
\begin{document}
\section*{UnsolvedMath ID: 700253}
\textbf{Source problem code:} catalogue entry
\subsection{Definitions and mathematical statement}
A group identity is an equation \(w(x_1,\ldots,x_n)=1\), where \(w\) is a group word, required to hold for every substitution of group elements for its variables. A variety of groups is the class of all groups satisfying some set of group identities. It is finitely based if it can be defined, among all groups, by finitely many identities. A subvariety of \(\mathcal V\) is a variety contained in \(\mathcal V\); it is proper if it differs from \(\mathcal V\).

A variety \(\mathcal V\) is a \emph{limit variety} if \(\mathcal V\) is not finitely based and every proper subvariety \(\mathcal W\subsetneq\mathcal V\) is finitely based. Here finite bases are bases among all groups, rather than bases relative to \(\mathcal V\). For a group \(G\), let \(\operatorname{var}(G)\) denote the smallest variety containing \(G\), equivalently the variety defined by all identities valid in \(G\).

\begin{enumerate}
\item[(a)] Give an explicit limit variety of groups, either by specifying a set \(\Sigma\) of identities whose class of models is a limit variety, or by specifying a group \(G\) such that \(\operatorname{var}(G)\) is a limit variety. The requested example must therefore fail to have a finite identity basis, while each of its proper subvarieties has such a basis.
\item[(b)] The archived question asks whether the collection \(\mathcal L\) of all limit varieties of groups is countable. The supplied Archive of Solved Problems entry reports a negative answer, with
\[
|\mathcal L|=2^{\aleph_0}.
\]
This is the supplied archived resolution, not a result independently verified in this draft.
\end{enumerate}

The source also states, as background attributed to Zorn's lemma, that every variety of groups having no finite basis of identities contains a limit subvariety. The supplied record labels the whole problem partially solved, retaining part (a) in the main section and placing part (b) in the solved archive.
\subsection{Short English statement}
A limit variety of groups has no finite identity basis, but every smaller variety inside it does. Part (a) asks for an explicit example, described by identities or a generating group. Part (b) asks whether there are only countably many such varieties; the supplied archive reports that there are continuum many. The record retains part (a) as unresolved.
\subsection{Sources}
\begin{itemize}
\item Local UnsolvedMath registry, permanent ID \(700253\), exact source code \texttt{catalogue entry}, titled \emph{Kourovka Notebook Problem 4.46}; supplied upstream revision \texttt{d371ee917ac9647fbb06cd6fc529deec6f3dda82}.
\item A.\ Yu.\ Olshanskii, proposer of Problem 4.46, \emph{The Kourovka Notebook (Unsolved Problems in Group Theory)}, 4th issue (1973). The supplied record describes an update in the 21st issue (2026), identified as \texttt{arXiv:1401.0300v48}. The collection description names E.\ I.\ Khukhro and V.\ D.\ Mazurov as editors.
\item P.\ A.\ Kozhevnikov, \emph{On varieties of groups of large odd exponent}, Dep.\ 1612-V00, VINITI, Moscow, 2000 (Russian). Cited in the supplied solved archive entry for the assertion that there are continuum many limit varieties.
\item S.\ V.\ Ivanov and A.\ M.\ Storozhev, \emph{Contemp. Math.}, volume 360 (2004), pp.\ 55--62. Cited in the same supplied archive entry; no article title is supplied.
\item The supplied status review is dated \texttt{2026-10-16} and labels the record \texttt{partially\_solved}, with archived solved part(s) and remaining part(s) in the main section. This assessment is supplied dated curation, not independently verified literature assessment. The full background supplies bibliographic details truncated in the status-review note.
\end{itemize}
\end{document}
